\section{Minmax Polynomials}

Weierstrass's theorem shows that polynomials approximate every
continuous function arbitrarily well if their degrees are
unrestricted. We now fix the degree and seek the smallest
possible uniform error.

\begin{definition}[Minmax Polynomial]
Let $f\in C([a, b])$ and $n\in\Nzero$. A polynomial
$p_n^*\in P_n$ is called a \emph{minmax polynomial} of $f$
of degree at most $n$ if
\[
 \|f-p_n^*\|_\infty
 \le\|f-p\|_\infty\qquad(p\in P_n).
\]
Equivalently, $p_n^*$ minimizes
$\max_{x\in[a, b]}|f(x)-p(x)|$ over $p\in P_n$.
\end{definition}

A minimum need not exist for an arbitrary optimization problem:
$x^2$ has no minimizer on $\{x\in\R\mid|x|>1\}$.
Nor must it be unique: on $\{x\in\R\mid|x|\ge1\}$,
its two minimizers are $-1$ and $1$. Existence and uniqueness
of minmax polynomials therefore require justification.

\Needspace{11\baselineskip}
\begin{example}[Degree-Zero Minmax Polynomial]
Let $m=\min_{x\in[a, b]}f(x)$ and
$M=\max_{x\in[a, b]}f(x)$. Since $f$ is continuous,
these extrema are attained. For any constant $c$,
\[
 \|f-c\|_\infty=\max\{|M-c|, |m-c|\}
 \ge\frac{M-m}{2}.
\]
The constant
\[
 c=\frac{M+m}{2}
\]
achieves equality, so it is a degree-zero minmax polynomial,
with error $(M-m)/2$. It is unique: if the error is at most
$(M-m)/2$, then necessarily
\[
 M-\frac{M-m}{2}\le c\le m+\frac{M-m}{2},
\]
forcing $c=(M+m)/2$.
\end{example}

\begin{theorem}[Existence of Minmax Polynomials]
Let $a<b$, $f\in C([a, b])$, and $n\in\Nzero$. Then there
exists a minmax polynomial $p_n^*\in P_n$ of $f$.
\end{theorem}
\begin{proof}
The vector space $P_n$, equipped with $\|\cdot\|_\infty$,
is finite-dimensional. Fix
\[
 C=2\|f\|_\infty+1>0,
 \qquad
 K=\{p\in P_n\mid\|p\|_\infty\le C\}.
\]
Since closed bounded subsets of a finite-dimensional normed
space are compact, $K$ is compact. Moreover,
$p\mapsto\|f-p\|_\infty$ is continuous because
\[
 \bigl|\|f-p\|_\infty-\|f-q\|_\infty\bigr|
 \le\|p-q\|_\infty.
\]
By the Extreme Value Theorem, some $p_n^*\in K$ minimizes
$\|f-p\|_\infty$ over $K$. Since $0\in K$,
\[
 \|f-p_n^*\|_\infty\le\|f\|_\infty.
\]
On the other hand, if $q\in P_n\setminus K$, then
\[
 \|f-q\|_\infty
 \ge\|q\|_\infty-\|f\|_\infty
 >C-\|f\|_\infty
 >\|f-p_n^*\|_\infty.
\]
Thus no polynomial outside $K$ improves the minimum, and
$p_n^*$ minimizes the error over all of $P_n$.
\end{proof}
